\chapter{La rotation de teinte}

\section{Un défaut que la mesure de chroma ne voyait pas}
Les premières versions du modèle reproduisaient le gain de chroma de
Lightroom à quelques centièmes près, et pourtant l'œil trouvait la vibrance
d'\IPR{} « moins forte ». Les couleurs ne basculaient pas vers leurs teintes
voisines : le bleu devait tendre vers le cyan. Une mesure qui ne regarde que la
chroma ne peut pas voir ce défaut, puisqu'une rotation conserve la chroma.

\resultat{Une métrique ne valide que ce qu'elle mesure. L'œil a révélé une
dimension, la teinte, que l'indicateur choisi ignorait.}

\section{D'où vient la rotation ?}
Augmenter la saturation dans un espace RGB revient à éloigner le pixel du gris
le long d'une droite de cet espace. Or les droites d'un espace RGB ne sont pas
des lignes de teinte constante pour la vision : en Oklab, leur teinte dérive.
C'est le même phénomène qui fait virer au violet les bleus saturés dans
d'autres espaces perceptuels, ou qui rend les oranges plus jaunes quand on les
pousse.

Les rotations mesurées sont cohérentes avec une vibrance calculée dans un
espace RGB de type ProPhoto : les bleus partent vers le cyan, les oranges vers
le jaune, les verts ne bougent presque pas. Cela reste une interprétation : ce
qui compte pour le modèle, c'est la rotation mesurée, quelle qu'en soit la
cause.

\section{La table de rotation}
La rotation est mesurée et stockée exactement comme les gains : pour chaque
valeur du curseur et chaque cellule (saturation, clarté, teinte), la médiane
de $\Delta h$ sur les scènes. Trois précautions la distinguent :
\begin{itemize}
\item $\Delta h$ est ramené à $]-180^\circ,180^\circ]$ avant toute médiane :
une teinte passant de $359^\circ$ à $1^\circ$ a tourné de $+2^\circ$, pas de
$-358^\circ$ ;
\item le repli des cellules vides conserve la teinte (chapitre 4) ;
\item entre les valeurs du curseur, l'angle est interpolé linéairement, en
degrés, avec zéro à zéro.
\end{itemize}

Dans le module, la rotation s'applique avec la matrice de rotation plane de
l'équation~\eqref{eq:transformation}, après le gain de chroma. Comme gain et
rotation commutent, l'ordre n'a pas d'importance.

\section{Ce que la rotation a corrigé}
Rotation médiane mesurée sur le rendu d'\IPR{} avant et après l'ajout de la
table, comparée à Lightroom, à $+100$ :

\begin{center}
\begin{tabular}{lrrr}
\toprule
Famille & Lightroom & \IPR{} avant & \IPR{} après\\
\midrule
rouge & $+1{,}35^\circ$ & $0^\circ$ & $+1{,}44^\circ$\\
orange & $+3{,}31^\circ$ & $0^\circ$ & $+3{,}30^\circ$\\
jaune-orangé & $+4{,}61^\circ$ & $0^\circ$ & $+4{,}82^\circ$\\
jaune & $+1{,}48^\circ$ & $0^\circ$ & $+1{,}58^\circ$\\
cyan & $-3{,}62^\circ$ & $0^\circ$ & $-4{,}67^\circ$\\
bleu-cyan & $-7{,}67^\circ$ & $0^\circ$ & $-4{,}55^\circ$\\
bleu & $-5{,}13^\circ$ & $0^\circ$ & $-5{,}10^\circ$\\
\bottomrule
\end{tabular}
\end{center}

À $+50$ et $-50$, l'écart de rotation reste inférieur à $0{,}6^\circ$ pour toutes
les familles. Le bleu-cyan à $+100$ reste en retrait ; le chapitre 6 montre
qu'il s'agit surtout d'un grand ciel bleu saturé, sur lequel Lightroom pousse
plus fort que son comportement médian.
